% =============================================================================
%  Section IV — Strategy Execution Models
% =============================================================================

\section{Strategy Execution Models}
\label{sec:execution-models}

Load-balancing agents in distributed edge systems must decide not only \emph{how}
to redistribute traffic — the strategy logic — but also \emph{when} to act:
at fixed intervals, in immediate response to peer events, or as a combination of
both.
Prior \ac{DFaaS} releases coupled these concerns: each strategy hosted its own
goroutine with a hardcoded \texttt{time.Ticker}.
While adequate for periodic recalculation, this design forced all future
strategies into the same mold, precluding the reactive behavior that becomes
natural once strategies have access to asynchronous peer signals such as
\texttt{MsgOverloadAlert} and \texttt{MsgFunctionEvent}.
This section describes the execution model abstraction introduced to decouple the
\emph{when} from the \emph{what}, enabling a strategy developer to declare an
execution model through Go's structural type system without modifying any
framework code.

% -----------------------------------------------------------------------
\subsection{Interface Hierarchy}
\label{subsec:interfaces}

The execution model is expressed through three composable Go interfaces.
\texttt{PeriodicStrategy} requires implementations to provide a
\texttt{Period()~time.Duration} method and a
\texttt{Tick(ctx~context.Context)~error} method; the framework invokes
\texttt{Tick} on every ticker event and threads the request context through so
that \texttt{Tick} bodies can respect cancellation during long operations such
as \ac{HAProxy} configuration pushes or inter-phase pauses.
\texttt{EventDrivenStrategy} additionally requires
\texttt{TriggerEvents()~[]string} — returning the message type constants that
should trigger execution — and \texttt{Debounce()~time.Duration}, which sets the
minimum interval between consecutive \texttt{React(ctx,~StrategyEvent)} calls;
a debounce of zero means \texttt{React} is called immediately on each trigger
arrival.
\texttt{HybridStrategy} embeds both constituent interfaces through Go's anonymous
composition; because both declare \texttt{OnReceived} with identical signatures,
the Go specification merges them into a single method requirement, so
implementors provide exactly one \texttt{OnReceived} body.

This three-level hierarchy applies the \emph{Strategy} design pattern~\cite{gamma1994patterns},
with the execution model selected at construction time rather than hardcoded in
each strategy.
All three interfaces include \texttt{OnReceived(msg~*pubsub.Message)~error}.
This method is invoked for every incoming broadcast message — including types
not listed in \texttt{TriggerEvents} — and is intended exclusively for state
table updates.
Separating state updates (universally called) from decision triggers (selective)
ensures that strategies maintain fresh peer state irrespective of their execution
cadence.
The \texttt{StrategyEvent} struct forwarded to \texttt{React} carries the type
discriminator and the full wire envelope as a \texttt{json.RawMessage}, allowing
\texttt{React} implementations to unmarshal into the concrete message type they
expect without an additional network round-trip.

% -----------------------------------------------------------------------
\subsection{Runner Implementations}
\label{subsec:runners}

The \texttt{StrategyRunner} interface abstracts strategy lifecycle from the rest
of the agent through two methods: \texttt{Callback()}, which returns the pubsub
callback to register with the \ac{GossipSub} receiver, and
\texttt{Run(ctx~context.Context)~error}, which blocks until the context is
cancelled or a fatal error occurs.
The \texttt{NewRunner} factory function examines the concrete type of the
strategy via a type-switch, placing \texttt{HybridStrategy} first to prevent it
from being spuriously matched as one of its constituent interfaces — a
requirement that arises because Go evaluates type-switch cases in declaration
order and \texttt{HybridStrategy} satisfies both \texttt{PeriodicStrategy} and
\texttt{EventDrivenStrategy}.

\paragraph{Periodic runner.}
The \texttt{periodicRunner} maintains a \texttt{time.Ticker} initialized from
\texttt{Period()} — defaulting to one minute when \texttt{Period()} returns zero
— and blocks in a \texttt{select} loop over the ticker channel and the context's
\texttt{Done} channel, invoking \texttt{Tick} on each tick event.
Errors from \texttt{Tick} are logged at WARN level and do not interrupt the
loop; transient failures in a single decision cycle — for example, a brief
\ac{HAProxy} \ac{API} timeout — should not terminate the agent.

\paragraph{Event-driven runner.}
The \texttt{eventDrivenRunner} introduces a capacity-one channel between the
callback goroutine and the worker goroutine.
A shared \texttt{makeTriggerCallback} helper builds the callback closure: on
each incoming message it calls \texttt{OnReceived}, then deserializes the
\texttt{MsgEnvelope} header and performs a set-membership check against a
pre-computed trigger set (built at construction from \texttt{TriggerEvents()}
using a hash-map for $\mathcal{O}(1)$ lookup); matching events are forwarded to
the channel with a non-blocking send, so excess events are dropped rather than
blocking the callback goroutine.
When debounce is configured, the worker does not invoke \texttt{React}
immediately; instead it records the latest event as \emph{pending} and resets a
\texttt{time.Timer} using Go's stop-then-drain-then-reset sequence to avoid
stale firings caused by the non-atomic semantics of
\texttt{time.Timer.Stop}.
When the timer fires, the most recent pending event is dispatched to
\texttt{React}, collapsing bursts of rapid arrivals into a single decision
cycle.
On context cancellation, any active debounce timer is explicitly stopped before
the goroutine returns, releasing its internal timer goroutine.

\paragraph{Hybrid runner.}
The \texttt{hybridRunner} unifies both execution models in a single
\texttt{select} statement over the ticker channel and the event channel.
\texttt{Tick} and \texttt{React} are therefore never concurrent by construction,
requiring no additional synchronization.
The event-delivery policy is channel-drop rather than debounce: if a trigger
event arrives while \texttt{Tick} or \texttt{React} is executing, it occupies
the capacity-one buffer; subsequent events are discarded until the buffer is
drained.
This policy bounds memory use and prevents the callback goroutine from blocking
on a slow worker, at the cost of potentially missing intermediate events — a
trade-off appropriate for strategies whose \texttt{React} logic re-reads current
system state rather than accumulating individual event payloads.

% -----------------------------------------------------------------------
\subsection{Context Propagation and Lifecycle}
\label{subsec:context}

All three runners propagate the agent's root context into every \texttt{Tick}
and \texttt{React} invocation, enabling strategy implementations to respect
cancellation during \ac{I/O}-bound operations.
The root context is created with \texttt{context.WithCancel} at agent startup;
its cancellation function is called when the process receives \texttt{SIGINT} or
\texttt{SIGTERM}, atomically signaling shutdown to all goroutines — including the
runner — through channel closure.
This design replaces earlier approaches that relied on process termination to
stop goroutines, enabling graceful shutdown of in-flight \ac{HAProxy}
configuration updates.
The \texttt{sleepOrCtx} helper, used by \texttt{RecalcStrategy} for the
mandatory inter-phase pause between its two-step recalculation cycle, implements
a context-aware \texttt{time.After} that returns the context error immediately if
cancellation arrives during the wait, rather than forcing the runner to block for
the full pause duration before observing shutdown.

% -----------------------------------------------------------------------
\subsection{Migration of Existing Strategies}
\label{subsec:migration}

All four production strategies — \texttt{RecalcStrategy},
\texttt{NodeMarginStrategy}, \texttt{StaticStrategy}, and
\texttt{AllLocalStrategy} — were migrated to the \texttt{PeriodicStrategy}
interface with no change to their internal decision logic.
Each migration followed the same mechanical transformation: the self-managed
goroutine and \texttt{time.Ticker} were removed; the ticker period was promoted
to a \texttt{Period()} method; and the loop body became the \texttt{Tick()}
method.
State that previously resided as locals inside the loop and needed to persist
across invocations was promoted to struct fields initialized in the factory's
\texttt{createStrategy()} method — most notably, the previous-tick function list
in \texttt{AllLocalStrategy}, which detects function set changes to trigger
incremental \ac{HAProxy} updates, and the resource threshold assignments and
initial node classification in \texttt{NodeMarginStrategy}.
Each migrated type carries a compile-time interface assertion
(\texttt{var~\_~PeriodicStrategy~=~(*XStrategy)(nil)}) to guard against future
accidental regressions caused by method-signature drift.
The migration is fully backward-compatible at the configuration level: the
\texttt{AGENT\_STRATEGY} environment variable continues to select among the same
strategy names, and the observable agent behavior is unchanged.
